\section{Konklusjon}
Designet vi har presentert viser hvordan bataljonsstridsgruppens ugraderte nettverk kan organiseres i adskilte sikkerhetssoner med kontrollert kommunikasjon mellom tjenestene og lokasjonene. Dette, sammen med tilgangskontroll, sikker administrasjon og logging samt gjennomføring av de foreslåtte fysiske og organisatoriske tiltakene, vil redusere risikoen knyttet til det flate nettverket.

Laboratorietestene viser at sentrale deler av designet fungerer etter hensikten. Dette gjelder segmentering og ruting, filtrering, kryptert transport, tilgangskontroll og logging. Derimot har ikke alle deler av designet blitt testet. ERSPAN er ikke verifisert ende til ende, og plattformspesifikk IOS-kompatibilitet må kontrolleres mot det endelige utstyret. Det må også vurderes om løsningen kan driftes innenfor dagens ressursramme.

Designet gir dermed S6 et grunnlag for videre verifikasjon og beslutning om rekonfigurasjon. Før endringer i produksjonsnettet bør de uavklarte forholdene testes på faktisk utstyr og transportforbindelse, det bør bekreftes at eksisterende tjenester opprettholdes, og gjenværende operativ risiko bør vurderes og aksepteres av S6.

